Целью данного раздела является доказательство того, что каждое облако
является собственным классом. Нам понадобится следующая вспомогательная лемма.

\begin{lemma}[\cite{levySet}]
  У любого множества кардинальных чисел есть верхняя грань.
\end{lemma}
\begin{corollary}
  Всякий класс кардиналов, не ограниченных сверху, является собственным.
  \label{colCardinal}
\end{corollary}

Далее, сформулируем и докажем теорему о классе пространств в каждом облаке.
\begin{theorem}
  Все облака представляют собой собственные классы.
\end{theorem}
\begin{proof} По следствию \ref{colCardinal} достаточно показать, что в любом
  облаке лежат пространства сколь угодно большой мощности.  Пусть $X$ --
  метрическое пространство мощности \(\alpha\), \(\Delta_\beta\)  ---
  симплекс мощности
  $\beta$, где $\beta > \alpha$. Обозначим $X_\beta = X \sqcup \Delta_\beta$.
  Зафиксируем произвольную точку $x$ пространства $X$ и положим
  расстояние от нее
  до любой точки симплекса равным $1$. Для точек $x' \in X$,
  $y \in \Delta_\beta$ определим
  \[\rho_{X_\beta}(y,x') = \rho_{X_\beta}(x',y) := \rho_X(x',x) + 1.\]
  Расстояния между другими парами точек оставим без изменений.
  Симметричность и неотрцательность расстояния $\rho_{X_{\beta}}$ очевидны.
  Для того чтобы
  полученное расстояние являлось метрикой достаточно проверить выполнение
  неравенства треугольника
  $\rho_{X_\beta}(x',z') \le \rho_{X_\beta}(x',y') +\rho_{X_\beta}(y',z')$
  только в том случае, если точки $x', y', z'$ не лежат одновременно в
  $\Delta_\beta$ или в $X$. Случаи $x', z' \in \Delta_\beta$ и $ x', z' \in X$
  очевидны. Разберем подробнее случаи, когда $x' \in X, z' \in \Delta_\beta$:
  \[ y' \in X: \rho_{X_\beta}(x', z') = \rho_X(x,x') + 1 \le
  \rho_X(x,y') + \rho_X(y',x') + 1 = \rho_X(x',y') + \rho_X(y',z'),\]
  \[y' \in \Delta_\beta: \rho_{X_\beta}(x', z') = \rho_X(x,x') + 1
  \le \rho_X(x',x) + 2 = \rho_X(x',y') + \rho_X(y',z').\]
  Итак, полученное пространство действительно будет метрическим. Осталось
  заметить, что если вложить $X$ в $X_\beta$, то $X_\beta$ будет лежать в
  замкнутой окрестности $X$ радиуса 1. Отсюда \( d_{GH}(X, X_ \beta )
  \le 1\), то есть \( X_ \beta  \) лежит в облаке \( [X] \).
\end{proof}

\begin{remark} Поскольку все облака являются собственными классами, между
  любыми двумя облаками существует биекция. Это означает, в
  частности, что класс соответствий между любыми двумя облаками не пуст.
\end{remark}
\begin{definition}
  Пусть $\mathcal{R}\big([X],[Y]\big)$ --- класс всех соответствий между
  облаками $[X]$ и $[Y]$. Определим \emph{искажение} соответствия $\dis R$
  аналогично определению \ref{def:Sootvet}.

  \emph{Расстоянием  Громова--Хаусдорфа} между облаками
  будем называть величину
  \[
    d_{GH}\big([X],[Y]\big) = \frac{1}{2}\inf\left\{\dis R : R\in
    \mathcal{R}\big([X],[Y]\big)\right\}.
  \]

\end{definition}
